% Generated draft: requires mathematical and source review.
\documentclass[11pt]{article}
\usepackage[margin=25mm]{geometry}
\usepackage{fontspec}
\setmainfont{Latin Modern Roman}
\usepackage{amsmath,amssymb,mathtools,unicode-math}
\setmathfont{Latin Modern Math}
\usepackage{enumitem,needspace,xurl,hyperref}
\setlength{\parindent}{0pt}
\setlength{\parskip}{5pt}
\setlength{\emergencystretch}{3em}
\newcommand{\N}{\mathbb{N}}
\newcommand{\Z}{\mathbb{Z}}
\newcommand{\Q}{\mathbb{Q}}
\newcommand{\R}{\mathbb{R}}
\newcommand{\C}{\mathbb{C}}
\newcommand{\F}{\mathbb{F}}
\newcommand{\A}{\mathbb{A}}
\newcommand{\PP}{\mathbb{P}}
\newcommand{\E}{\mathbb{E}}
\newcommand{\eps}{\varepsilon}
\newcommand{\dd}{\,\mathrm{d}}
\begin{document}
\section*{UnsolvedMath ID: 800629}
\textbf{Source problem code:} catalogue entry
\subsection{Definitions and mathematical statement}
The supplied record leaves the class of spaces and the cover conventions unspecified. The following is a conditional formulation using standard topological meanings of the notation; the intended class and conventions require source verification.

Fix a class \(\mathcal C\) of topological spaces. For each \(X\in\mathcal C\), fix a collection \(\mathscr O(X)\) of admissible open covers of \(X\), with \(X\notin\mathcal U\) for every \(\mathcal U\in\mathscr O(X)\). An open cover is a family of open subsets whose union is \(X\). Whether admissible covers must be countable is unresolved in the supplied record. Define
\[
\Omega(X)=\left\{\mathcal U\in\mathscr O(X): (\forall F\subseteq X\text{ finite})(\exists U\in\mathcal U)\ F\subseteq U\right\}
\]
and
\[
\Gamma(X)=\left\{\mathcal U\in\mathscr O(X):\mathcal U\text{ is infinite and }(\forall x\in X)\ \{U\in\mathcal U:x\notin U\}\text{ is finite}\right\}.
\]
For classes \(\mathcal A(X),\mathcal B(X)\) of covers, the property \(\mathsf S_{\mathrm{fin}}(\mathcal A,\mathcal B)\) means that for every sequence \((\mathcal U_n)_{n\geq1}\) in \(\mathcal A(X)\), there are finite subfamilies \(\mathcal F_n\subseteq\mathcal U_n\) such that
\[
\bigcup_{n\geq1}\mathcal F_n\in\mathcal B(X).
\]
The property \(\mathsf U_{\mathrm{fin}}(\mathcal A,\mathcal B)\) means that for every sequence \((\mathcal U_n)_{n\geq1}\) in \(\mathcal A(X)\), each \(\mathcal U_n\) having no finite subcover of \(X\), there are finite subfamilies \(\mathcal F_n\subseteq\mathcal U_n\) such that
\[
\left\{\bigcup\mathcal F_n:n\geq1\right\}\in\mathcal B(X).
\]
Here finite selections may be empty, and \(\bigcup\varnothing=\varnothing\); this convention also requires verification against the original source.

With the intended class \(\mathcal C\) and cover conventions fixed, Problem 1 asks whether
\[
(\forall X\in\mathcal C)\quad
X\models\mathsf U_{\mathrm{fin}}(\Gamma,\Omega)
\ \Longleftrightarrow\
X\models\mathsf S_{\mathrm{fin}}(\Gamma,\Omega).
\]
If this equivalence fails, it further asks whether
\[
(\forall X\in\mathcal C)\quad
X\models\mathsf U_{\mathrm{fin}}(\Gamma,\Gamma)
\ \Longrightarrow\
X\models\mathsf S_{\mathrm{fin}}(\Gamma,\Omega).
\]
The unspecified class of spaces and cover conventions prevent identifying this conditional formulation with a uniquely determined source statement. The registry labels the problem open; this is supplied status information, not an independently verified mathematical or literature assessment.
\subsection{Short English statement}
For the intended spaces and cover conventions, ask whether finite selections from sequences of \(\gamma\)-covers give equivalent properties when their individual members or their unions must form an \(\omega\)-cover. If the properties differ, ask whether the corresponding union-selection property with a \(\gamma\)-cover as output implies \(\mathsf S_{\mathrm{fin}}(\Gamma,\Omega)\). The intended domain and conventions remain to be verified.
\subsection{Sources}
\begin{itemize}
\item Local UnsolvedMath registry, supplied complete source record: permanent ID \(800629\), exact source problem code \texttt{catalogue entry}, title ``Selection Principles Milestone Problems --- Problem 1.'' The record attributes the problem to Tsaban, with proposed year 2003, and supplies status ``open.''
\item Supplied source list: Tsaban, \emph{Selection Principles in Mathematics: A Milestone of Open Problems} (2003), Problem 1, \url{https://arxiv.org/abs/math/0312182}. The record reports extraction from arXiv LaTeX and access on 2026-10-16; these are supplied provenance details, not independently verified access or a current literature assessment.
\item The supplied statement includes the reference label ``coc2, Problems 1 and 2.'' The corresponding bibliography entry is absent from the supplied record, so its authorship and publication details remain unresolved.
\end{itemize}
\end{document}
